\DocumentMetadata{
  pdfversion=2.0,
  pdfstandard=ua-2,
  lang=en-US,
  tagging=on,
  tagging-setup={math/setup=mathml-SE,tabsorder=structure}
}
\documentclass[11pt,letterpaper]{article}
\usepackage[margin=0.68in,includefoot]{geometry}
\usepackage{newcomputermodern}
\usepackage{amsmath,amscd,mathtools}
\usepackage{tikz,adjustbox}
\providecommand{\square}{\mdlgwhtsquare}
\usepackage[hidelinks]{hyperref}
\hypersetup{
  pdftitle={Algebra First-Year Exam, August 2021, Part 1},
  pdfauthor={University of Florida Department of Mathematics},
  pdfsubject={Graduate mathematics examination},
  pdfdisplaydoctitle=true,
  bookmarksopen=true
}
\setcounter{secnumdepth}{0}
\setlength{\parindent}{0pt}
\setlength{\parskip}{0.28em}
\setlength{\emergencystretch}{3em}
\setlength{\leftmargini}{2.15em}
\setlength{\leftmarginii}{2.65em}
\setlength{\leftmarginiii}{3em}
\renewcommand{\labelenumi}{\textbf{\theenumi.}}
\pagestyle{empty}
\raggedbottom
\allowdisplaybreaks
\newenvironment{examproblems}[1]{%
  \begin{enumerate}%
  \setcounter{enumi}{#1}%
  \setlength{\topsep}{0.35em}%
  \setlength{\partopsep}{0pt}%
  \setlength{\itemsep}{0.48em}%
  \setlength{\parsep}{0.18em}%
}{\end{enumerate}}
\newenvironment{examsubparts}{%
  \begin{enumerate}%
  \setlength{\topsep}{0.18em}%
  \setlength{\partopsep}{0pt}%
  \setlength{\itemsep}{0.16em}%
  \setlength{\parsep}{0.1em}%
}{\end{enumerate}}
\newenvironment{examinstructions}{%
  \par\begingroup\itshape
}{%
  \par\endgroup\vspace{0.18em}
}
\makeatletter
\renewcommand\section{\@startsection{section}{1}{0pt}%
  {-1.25ex plus -.25ex minus -.1ex}%
  {0.55ex plus .1ex}%
  {\normalfont\large\bfseries}}
\renewcommand{\@maketitle}{%
  \newpage\null\vskip -1em%
  {\footnotesize\bfseries
    \href{https://www.ufl.edu/}{University of Florida}, \href{https://math.ufl.edu/}{Department of Mathematics}\par}%
  \vskip -0.5em%
  \tagstructbegin{tag=Title}%
    {\LARGE\bfseries\href{https://gma.math.ufl.edu/exams/algebra-first-year/}{Algebra First-Year Exam}, August 2021, Part 1\par}%
  \tagstructend%
  \vskip 0.25em%
  \tagmcbegin{artifact}%
    \hrule height 1pt%
  \tagmcend%
  \vskip 1em%
}
\makeatother
\title{Algebra First-Year Exam, August 2021, Part 1}
\author{}
\date{}
\begin{document}
\maketitle
\thispagestyle{empty}
\begin{examinstructions}
Answer four problems. Write your answers clearly in complete English sentences. You may quote results (within reason) as long as you state them clearly.
\end{examinstructions}
\begin{examproblems}{0}
\item
Let~\(G\) be a cyclic group. Prove that every subgroup of~\(G\) is cyclic. (Be sure to consider infinite cyclic groups as well as finite ones.)
\item
Let \(G, H\) be finite groups with orders \(m, n\).
\begin{examsubparts}
\item[\textnormal{(a)}]
Prove that if~\(m\) and~\(n\) are coprime then
\[
\operatorname{Aut}(G \times H) \cong \operatorname{Aut}(G) \times \operatorname{Aut}(H).
\]
\item[\textnormal{(b)}]
Give an example where~\(m\) and~\(n\) are not relatively prime and
\[
\operatorname{Aut}(G \times H) \not\cong \operatorname{Aut}(G) \times \operatorname{Aut}(H).
\]
 Justify your claims.
\end{examsubparts}
\item
This problem has two parts.
\begin{examsubparts}
\item[\textnormal{(a)}]
Let~\(G\) be a finite group of order~\(n\) and let~\(p\) be the smallest prime which divides~\(n\). Prove that if~\(H\) is a subgroup of~\(G\) of index~\(p\) then~\(H\) is normal in~\(G\).
\item[\textnormal{(b)}]
Give an example of a finite group~\(G\) and a subgroup \(H \leq G\) such that~\(H\) is not normal in~\(G\) and the index \(|G:H|\) is prime.
\end{examsubparts}
\item
Let \(p,q,r\) be primes such that \(p<q<r\) and let~\(G\) be a group of order \(pqr\). Prove that~\(G\) has a normal Sylow subgroup.
\item
Prove that the alternating group \(A_5\) is simple. Prove that \(A_4\) is solvable.
\item
Show that there are four different homomorphisms \(\phi:\mathbb{Z}_2\to\operatorname{Aut}(\mathbb{Z}_8)\), and prove that the corresponding semidirect products \(\mathbb{Z}_8\rtimes_\phi\mathbb{Z}_2\) are pairwise nonisomorphic. (Hint: Consider the number of elements of order 2.)
\end{examproblems}
\end{document}
